\chapter{Decision}
\label{chap:decision}

\begin{chapteressay}
The end of the procession is not \lean{TRUTH}. It is \lean{INFERRED}. That is
not a retreat. It is the mature shape of public science. A decision made under
a record can be strong enough to build bridges, approve treatments, announce a
particle, revise cosmology, or regulate a risk. It still remains a decision
under a record.

Decision follows \lean{LOGICAL}, \lean{HALTED}, \lean{MEASURED},
\lean{COMPILED}, and \lean{INFERRED}. The question is no longer whether marks
can enter, residues can persist, values can be staged, overhead can be named,
or truth can travel. The question is what conclusion the whole chain licenses.

Decision is the most dangerous movement because it is where people are most
tempted to pretend the record has become command. A protocol shows an effect;
therefore a treatment must be used. A model compiles records; therefore the
model owns reality. A threshold is crossed; therefore doubt has become
disloyal. A climate chain licenses factual inference; therefore every policy
preference can borrow the authority of physics. Decision refuses all of
that. It also refuses the opposite laziness, where every finite record is
treated as too uncertain to matter. The record can bind without becoming
absolute.

The words in this movement are deliberately cooler than the earlier ones.
\lean{LOGICAL} is not the whole of reason. It is a disciplined relation among
premises, rules, and admissible conclusions. \lean{HALTED} is not the death of
inquiry. It is a public stopping condition for a particular decision under a
particular record. \lean{MEASURED} is not all that exists. It is what has
entered the record through instruments and protocols. \lean{COMPILED} is not
a magic fusion of evidence into certainty. It is the assembly of heterogeneous
records into a working structure. \lean{INFERRED} is the mature word because
it remembers that the conclusion has a route.

Decision therefore has a double ethic. It asks decision to be brave enough
to act under good records and humble enough to remember the record's
conditions. Bridges are not built from metaphysical certainty. Treatments are
not approved because a population average is identical to a patient. Particle
discoveries are not announced because all possible doubt has died. Climate
inferences are not factual because policy requires them, nor optional because
policy is contested. In each case the question is narrower and harder: what
does this record license?

That question also protects the earlier chapters from becoming romantic. A
mark that never reaches decision can remain a cabinet curiosity. A residue
that never informs inference can remain an irritation. A staged value that
never compiles can remain a local success. Transported truth that never enters
conclusion can remain admirable but inert. Decision is where the record is
asked to bear weight. The book does not avoid that weight. It asks how much
the record can bear without breaking.

Decision also changes the meaning of disagreement. Before the halt condition,
disagreement may be part of building the record. After a public inference,
disagreement may be about the record's scope, the values attached to action,
the adequacy of the model, the population to which a result applies, or the
cost of being wrong. A serious culture keeps those disagreements sorted. It
does not call every disagreement denial, and it does not call every denial a
brave disagreement. Sorting is part of decision.

The examples range from formal science to ordinary public judgment
because inference is not only a laboratory event. Double-blind trials show
desire being blocked before measurement can advise action. The Standard Model
shows many records becoming one durable structure without becoming final
ontology. Cosmology shows the universe becoming inferable through a chain
whose rungs differ in strength. Discovery thresholds show how a community
names public halt conditions under uncertainty. Climate compilation shows
factual inference and policy choice being kept in the same room without being
collapsed. Noether and Galois show formal closure as a seed of disciplined
inference. Juries, bridges, and treatments show how public life acts under
records that are strong, finite, and consequential.

The movement also answers a temptation left by the device discussion. If one
asks whether a final object called \(K\) exists, or whether the universe is a
simulation, experiment cannot simply compute the metaphysical proposition by
looking harder. What experiment can do is build a procession of public
records: marks, residues, staged values, overhead, transported witness, and
licensed inference. That procession may make some conclusions responsible and
others irresponsible. It does not turn inference into possession of the
universe from outside.

The book ends in \lean{INFERRED} because the word is not small because science
is weak. It is small because strong records do not need inflated
words. An inference can be the reason a bridge remains open or closes, a drug
is recommended or refused, a particle is named, a model is trusted for a
range, or a risk is regulated. That is enough seriousness. The conclusion
stays answerable because the route stays visible.

The route is also what lets this chapter handle error without theatrical
collapse. Decisions can be revised. A drug can be withdrawn, a bridge rating
changed, a particle anomaly dissolved, a model parameter updated, a climate
projection narrowed, a verdict appealed or historically judged. Revision does
not mean the earlier decision was irrational. It means the earlier decision
was made under a record, and the record later changed. A mature decision
culture knows how to say that without humiliation.

This is Decision's realism about time. Decisions happen before the future
has supplied all its evidence. Waiting is itself a decision when bridges age,
patients suffer, emissions continue, detectors run, or public language needs
a name for a discovered particle. The question is not how to avoid decision.
The question is how to make the decision carry its record, scope, uncertainty,
and revision conditions.

The final movement therefore has a different tempo from the previous ones.
Earlier chapters often slowed down to show how a record was born. This one
slows down at the point where a record is used. Use is a severe test of truth.
It exposes overclaim, underclaim, hidden values, bad thresholds, and missing
domains. It also reveals the dignity of a record that can bear real weight
without pretending to be more than it is.

The reader should expect fewer miracles here. Decision is not a final blaze
of certainty. It is a craft of admissible conclusion. The craft asks what has
been measured, what has been protected by protocol, what has been compiled,
where the halt condition lies, and which values enter after measurement. A
conclusion that answers those questions may be quiet, but it can still change
the world.

That quietness is the point. Decision does not need to shout because
the previous chapters have done the work of making records durable. Decision
is where durability becomes responsibility.

Responsibility is the last gauge.
It asks whether a finite public record can carry finite public weight.
The answer is yes, but only with its route attached.
\end{chapteressay}

\section{Double-Blind Trials: Protocol Against Desire}
\episodecoord{1900s onward | Decision :: LOGICAL -> MEASURED}

A trial is a machine for denying appetite control over inference. Blinding does
not make investigators morally pure. It changes the record so desire has fewer
places to write itself as measurement.

The desire problem is not a moral insult. It is an experimental fact about
human beings. Patients want to improve. Physicians want to help. Sponsors want
effects. Investigators want clean results. Observers notice what they expect
to notice. Participants report differently when they believe they have
received help or harm. None of this requires fraud. It is enough that
expectation has routes into the record. A trial protocol is logical when it
blocks those routes before they can masquerade as measurement.

Randomization is the first discipline. It refuses to let assignment follow
the preferences of patient, physician, investigator, or institution. The
point is not that randomization creates a perfect miniature of the world. It
creates a comparison whose known and unknown influences have fewer systematic
paths into one arm rather than another. The treatment effect becomes more
measurable because allocation has been denied a story.

Blinding adds another discipline. If a patient knows the assignment, reporting
can shift. If a clinician knows the assignment, encouragement, follow-up,
co-interventions, or classification can shift. If an outcome assessor knows
the assignment, ambiguous endpoints can tilt. Blinding does not eliminate all
bias. It removes some of the easiest routes by which desire and expectation
enter the value. That removal is a decision technology.

The word "technology" is important. A double-blind trial is not merely a
philosophy of caution. It is a constructed procedure: placebo or comparator,
allocation concealment, coded materials, endpoint definitions, adverse event
rules, data collection forms, analysis plans, stopping rules, and unblinding
conditions. The trial makes inference possible by building a local world in
which preference is harder to smuggle into measurement.

Trials belong in Decision rather than only in Staging because staging showed
how protocols produce values. Decision asks whether a value can advise
action. A clinical trial must cross from measured outcomes to treatment
judgment. That crossing is never automatic. The protocol can show an average
effect under specified conditions; the decision must ask whether the effect
is large enough, safe enough, relevant enough, and applicable enough to the
population or patient at hand.

The endpoint is one of the clearest examples of a halt condition.
Before a trial begins, investigators have to say what will count: mortality,
symptom score, viral load, tumor response, blood pressure, hospitalization,
quality of life, or some composite. If the endpoint is moved after seeing the
data, desire has reentered under the name of interpretation. Registration and
pre-specified analysis matter because they preserve the logical
shape of the decision before the result tempts everyone to redraw it.

The placebo is also philosophically richer than it looks. It is not a fake
thing in a simple sense. It is a control for expectation, encounter, ritual,
and background care. A placebo arm says: let the treatment effect answer not
against nothing, but against a constructed comparison that carries many of
the same human circumstances. The decision becomes more realistic because the
human setting is admitted rather than denied.

There is a danger in overpraising the double-blind form. Some interventions
cannot be blinded cleanly. Surgery, behavioral treatment, public health
measures, device interventions, and obvious side effects can break the form.
The lesson is not that unblinded evidence is useless. The lesson is that
every design must name the routes by which desire can enter. If blinding is
impossible, other guards must carry more weight: objective endpoints,
independent assessment, allocation concealment, replication, pragmatic
comparison, sensitivity analysis, or transparent uncertainty.

The population problem is just as important. A trial result is not a verdict
about every patient. Inclusion criteria, exclusion criteria, adherence,
background care, demographics, comorbidities, follow-up, and endpoint choice
all shape the domain of inference. A trial may be internally strong and still
travel poorly to another population. Truth Transport's rule returns inside
Decision: decision must know where the trial can travel.

Clinical inference is a good opening for Decision because it is close to
action, and action exposes the limits of knowledge. A statistically
significant benefit can be clinically small. A meaningful average benefit can
hide subgroup harm. A surrogate endpoint can fail to predict what patients
actually care about. A short trial can miss long-term risk. A clean protocol
can be hard to implement in ordinary care. None of these facts abolishes
trial evidence. They keep decision honest about what kind of inference has
been licensed.

The double-blind trial also carries an ethical symmetry. It protects the
record from desire, but it also protects participants from being used as
props in a foregone conclusion. If the answer were already known, the trial
would be unnecessary or unethical. If the answer is not known, the protocol
owes participants a route by which their exposure can produce public
knowledge. The logic is not cold. It is one way respect becomes procedural.

The modern trial therefore sits between care and knowledge. It is not pure
care, because some participants receive control or comparator treatment. It
is not pure knowledge, because people are being treated and harmed or helped.
The decision structure has to carry both obligations. A trial that forgets
care becomes exploitation. A trial that forgets knowledge becomes anecdotal
medicine. The protocol is the narrow bridge between them.

When the trial ends, the decision does not end. The record enters review,
replication, guideline development, post-market surveillance, comparative
effectiveness, and patient-level judgment. The first inference may license
use under conditions; later records may widen, narrow, or reverse that use.
\lean{HALTED} does not mean inquiry has died. It means a public decision point
has been reached under the record now available.

A trial also teaches the difference between evidence of effect and evidence
of decision-worthiness. An intervention may beat placebo while still being
too toxic, too expensive, too small in effect, too difficult to deliver, or
too poorly matched to the patients who most need help. Conversely, a modest
effect may matter greatly if the condition is severe, alternatives are poor,
and risks are acceptable. The measured contrast is the beginning of decision,
not the whole of it.

This is where adverse events become part of inference rather than afterthought
or public-relations nuisance. A treatment effect that improves one endpoint
while worsening another must be judged as a balance, not as a victory lap.
Harms may be rarer than benefits and therefore harder to detect. Long-term
harms may not appear during a short trial. Subgroup harms may hide inside a
favorable mean. A decision protocol that wants to be realistic must ask how
harm enters the record.

Missing data are another test of honesty. People drop out, miss visits, stop
treatment, cross over, die, or become impossible to follow. The missingness is
rarely innocent. A participant may leave because treatment failed, because it
worked, because side effects were intolerable, because logistics failed, or
because the disease changed. The analysis must decide what to do with absence.
That decision can move the conclusion. A guarded record therefore has to make
missingness visible.

Intention-to-treat analysis is one way trial logic protects decision. It keeps
participants in the groups to which they were assigned, preserving the
comparison created by randomization. Per-protocol analysis asks a different
question about those who followed the treatment. Both can be useful, but they
do not answer the same decision. The point is not to canonize one statistical
rule. The point is to know which decision each rule licenses.

Early stopping is equally delicate. A trial stopped early for apparent benefit
may save participants from inferior treatment, but it may also exaggerate
effect size if the early signal was unusually favorable. A trial stopped for
harm may protect participants while leaving some uncertainty unresolved. Data
monitoring committees exist because decision can enter before the planned end
of measurement. Even stopping is a decision under a record.

Multiplicity is one of desire's subtle doors. If investigators measure enough
outcomes, subgroups, time points, and transformations, some comparison may
look impressive by chance. Pre-specification disciplines that field of
possible stories. It does not forbid exploration; it labels exploration.
Exploratory findings can guide future trials, but they should not borrow the
authority of the primary decision unless the route warrants it.

The trial therefore embodies a grammar of inference: assign before knowing,
blind before observing, define before measuring, analyze under rule, disclose
uncertainty, and decide within scope. This grammar is \lean{LOGICAL} in the
practical sense. It does not float above the clinic. It organizes
the clinic so that measurement can advise action.

The patient's role should not disappear inside this machinery. Participants
are not merely data sources. Consent, monitoring, rescue rules, stopping
conditions, and follow-up are part of the ethical route. A trial that produces
clean numbers by treating people as inert material has failed Decision even
if the statistical contrast is legible. The record that licenses action must
also be a record produced under acceptable care.

This makes double-blind trials a useful answer to romantic medicine and
mechanical medicine at once. Romantic medicine trusts the healing encounter
so much that it risks mistaking expectation for effect. Mechanical medicine
trusts the aggregate so much that it risks forgetting the person. A good trial
disciplines the encounter, then a good treatment decision returns the
aggregate to a person without forgetting the discipline.

The history of trial registration adds another layer of realism. A published
paper is not the whole record if unpublished trials, changed endpoints, or
selective reporting shape what the public sees. Registration and reporting
norms make the invisible comparison set more visible. They ask not only what
the successful paper says, but what trials were begun, what outcomes were
planned, and what results did or did not appear. Decision needs that wider
record.

The trial also shows why statistical significance and practical significance
must remain separate. A large sample can make a tiny effect statistically
detectable. A small sample can miss a clinically important effect. The
decision has to ask about magnitude, uncertainty, baseline risk, competing
risks, patient values, feasibility, and alternatives. Statistics helps
decision; it does not replace judgment about importance.

When trials conflict, the compiled record becomes necessary. One guarded
record may be underpowered, narrow, flawed, or local. Systematic review and
meta-analysis can compile multiple trials, but they introduce their own
transport questions: comparability of populations, doses, endpoints, follow-
up, bias, publication, and analysis. A compiled treatment inference is
stronger only if the compilation preserves these differences.

The opening anchor does not end with a single triumphant trial. It ends with a
decision ecology. Randomized double-blind design is a central
device, but clinical decision also needs registration, replication,
surveillance, synthesis, and patient judgment. The protocol blocks desire at
the start; later institutions must keep desire from reentering as selective
memory, marketing, or overgeneralization.

\begin{phenom}{The Double-Blind Decision Effect~\cite{fisher1935}}
\label{ph:double-blind-decision}

\PhStatement
A decision protocol is logical when it blocks known routes by which desire can
enter the measurement.

\PhOrigin
Modern experimental design and randomized controlled trials developed tools
for separating treatment effect from expectation, allocation bias, and observer
choice.

\PhObservation
The record includes assignment, concealment, endpoint, analysis plan, and
measured outcome.

\PhConstraint
An inference about effect must preserve the protocol that protected the record
from the investigator's preferred conclusion.

\PhConsequence
Clinical inference is not final truth about a patient or population. It is
licensed action under a guarded record.
\end{phenom}

The decision effect is strongest when its modesty is preserved. A double-blind
trial does not make the investigator virtuous, the population universal, or
the future known. It makes one route of inference more trustworthy by denying
desire some of its easiest entrances. That is enough to matter. It can change
practice, save lives, stop ineffective treatment, or reveal harm.

The section also gives Decision its first general rule: a decision protocol
must name the vulnerability it is built to block. Blinding blocks expectation.
Randomization blocks selection by preference. Endpoint specification blocks
after-the-fact success. Concealment blocks anticipatory assignment. Analysis
plans block flexible storytelling. Each device is logical because it closes a
known route from desire to conclusion.

This rule will recur in particle physics, cosmology, climate, engineering,
and ordinary judgment. A community may use different tools, but the shape is
the same. Decision becomes public when the record is protected against the
ways people most want it to speak.

The section's last caution is that protection is never total. A blinded trial
can still have poor endpoints, inadequate follow-up, broken masking,
unrepresentative patients, sloppy conduct, missing data, or selective
publication. The protocol is not a halo. It is a set of guards. Decision
requires asking whether those guards were appropriate to the vulnerability
and whether the remaining uncertainty changes the action.

Broken masking is especially instructive. Side effects, laboratory changes,
or obvious clinical responses can reveal assignment even when the protocol
was formally blinded. Once assignment is guessed, expectation may reenter.
The trial has to ask whether the blind was credible, whether outcomes were
objective enough to withstand partial unblinding, and whether assessment
procedures protected the comparison. The name "double-blind" alone is not the
evidence. The functioning blind is.

Comparator choice is another decision hidden inside design. A placebo may be
appropriate when no established treatment exists or when the question truly
requires it. An active comparator may be ethically required when withholding
treatment would be unacceptable. A noninferiority trial asks yet another
question: whether a new treatment is not unacceptably worse than an existing
one under specified margins. Each comparator licenses a different conclusion.
The trial's logic begins before the first participant is enrolled.

Margins are decision statements too. What counts as clinically meaningful?
What harm is tolerable? What loss of efficacy is acceptable for a gain in
safety, convenience, cost, or access? These are not purely statistical
questions. They are values made explicit enough for measurement to answer.
Decision becomes dishonest when the value-laden margin is hidden behind the
technical analysis.

Subgroups add another layer of discipline. Patients differ. Age, sex,
genetics, disease severity, coexisting illness, baseline risk, and concurrent
treatment can change benefit and harm. But subgroup analysis is fertile
ground for false discovery if every slice of the population is searched after
the fact. A trial must distinguish pre-specified biological or clinical
questions from exploratory pattern hunting. The patient may be particular,
but the route to particularity still needs rules.

Pragmatic trials and explanatory trials also answer different decisions. An
explanatory trial may ask whether an intervention works under controlled
conditions. A pragmatic trial may ask whether it works in ordinary care, with
messier adherence, varied clinicians, and real delivery systems. Neither is
automatically superior. Each transports truth differently. Decision needs to
know whether it is borrowing evidence about ideal effect or usable effect.

The double-blind anchor therefore gives the whole chapter a working doctrine:
decision is not the moment when measurement stops being technical and becomes
mere judgment. Decision is the moment when technical form, uncertainty, value,
and consequence must be kept in one accountable sentence. A trial that can
write that sentence has earned public force.

The accountable sentence often has the form: in this population, under this
comparison, with this endpoint, over this follow-up, the intervention changed
this outcome by this amount with this uncertainty and these observed harms.
Everything important is in the qualifiers. Remove them, and the sentence
becomes advertising. Preserve them, and the sentence can guide care, review,
and further study.

This is also why trial reports need tables that look boring. Baseline
characteristics, flow diagrams, confidence intervals, adverse event tables,
protocol deviations, and follow-up counts are not bureaucratic clutter. They
are the public route by which the reader can decide whether the inference is
admissible. A glossy conclusion without these supports is not a trial result
in the full Decision sense.

The trial section therefore names a general public obligation: the stronger
the action invited by a record, the more visible the record's guards must be.
A small exploratory claim may survive with modest protection. A practice-
changing claim needs stronger protocol, clearer endpoints, better accounting
for harms, and more careful transport to the population in question. Decision
weight demands decision custody.

The same obligation applies to negative decisions. A trial can stop a harmful
or ineffective treatment from continuing as custom. That kind of decision is
less glamorous than approval, but it is one of medicine's most important
forms of care. The record does not only authorize doing. It can authorize
not doing, stopping, warning, or narrowing use. \lean{HALTED} can protect
patients by halting enthusiasm.

Negative results also test publication ethics. If only positive trials travel,
the public record becomes biased toward intervention. A treatment may look
better than it is because failed or harmful attempts disappeared. Trial
registries, reporting requirements, and systematic reviews are attempts to
make the negative decision record visible. Decision requires memory of what
did not work.

This is the first chapter where non-action becomes a full scientific outcome.
Earlier movements emphasized how marks, residues, and transported truths come
to public life. Decision adds that a licensed refusal can be just as real as a
licensed adoption. The protocol may tell us not to treat, not to expand, not
to announce, not to cross a load limit. Such refusals are not weakness. They
are records bearing consequence.

The trial anchor also makes "evidence hierarchy" more concrete. A hierarchy
is not a caste system for papers. It is a memory of vulnerability. Anecdote is
vulnerable to selection, memory, and expectation. Observational comparison is
vulnerable to confounding. Randomized blinded trials block some of those
routes but introduce others. Meta-analysis can increase power while importing
publication bias and heterogeneity. The decision maker must know the
vulnerability each level carries.

This prevents a crude worship of randomized trials. Some questions cannot be
randomized ethically or practically. Some rare harms appear only after wide
use. Some public health decisions depend on observational records. The double
blind trial is an anchor because it shows desire being blocked with unusual
clarity, not because every real decision has the same form.

The anchor therefore teaches transfer rather than imitation. Other fields
need not copy the placebo-controlled trial. They need to copy its honesty
about vulnerability. What route might make the conclusion false while leaving
the record apparently clean? What design blocks that route? What uncertainty
remains? These questions are portable even when the clinical form is not.

The section can begin Decision because it gives the movement its moral and
logical posture: do not trust yourself too much; build a record that does
not need your purity.

That posture is not cynicism about investigators. It is mercy toward them.
The protocol carries burdens that no individual intention should be asked to
carry alone.

It also gives critics a fair object. They can criticize the assignment,
endpoint, blind, analysis, or applicability rather than guessing at motives.
That is how disagreement becomes useful.

Useful disagreement is one of the trial's public products.

\section{Standard Model Compilation}
\episodecoord{1900s onward | Decision :: COMPILED -> INFERRED}

The Standard Model is not one experiment. It is a compilation: scattering
records, spectra, accelerator events, conservation laws, symmetries, failures,
and confirmed predictions brought into one working structure. A realistic
account must not pretend the compilation is simpler than it is.

This is the section where \lean{COMPILED} earns its name. A compiled model is
not a heap of results. It is a structure in which different records constrain
one another. A scattering cross section limits one possibility. A decay
channel opens another. A forbidden process remains absent. A symmetry
organizes families. A precision residue demands correction. A particle
discovery fills a predicted role. A null search closes a path. The model is
not in any one of these records. It is the disciplined pattern they make
together.

The danger is retrospective neatness. Once a model has become standard, its
history can be told as if all the pieces were walking toward their assigned
places. That is not how the record felt while it was being made. The
compilation passed through cloud chambers, bubble chambers, counter
experiments, accelerators, electroweak theory, quark models, gauge theory,
renormalization, neutral currents, deep inelastic scattering, jets, precision
electroweak tests, and the long search for the Higgs boson. The finished name
compresses a long argumentative ecology.

The ecology matters because compilation is not the same as consensus. A
consensus can form socially around a weak structure. A compilation becomes
strong when independent records become difficult to move independently. If
one changes a coupling, a mass, a symmetry assignment, or a particle content,
other records begin to complain. The structure earns authority by making the
record interdependent.

This interdependence is why the Standard Model is both powerful and limited.
It organizes an enormous range of particle phenomena with extraordinary
precision. It also leaves neutrino masses, dark matter, gravity, cosmological
questions, and deeper hierarchy problems outside its completed form. A
compiled structure can be the best working account in its domain without
becoming final metaphysics. That is exactly the humility Decision needs.

The Standard Model also shows how theory can be measured without being seen.
No detector photographs "the Standard Model." Detectors record tracks,
energies, momenta, hits, missing transverse energy, invariant masses,
branching ratios, and event populations. The model enters as the structure
that makes those records cohere and that predicts where new records should or
should not appear. The inference is not direct vision. It is compiled
constraint.

Detector language is important because it keeps the model from becoming a
floating diagram. A collision produces sprays, hits, deposits, timing,
curvature, missing energy, and reconstructed candidates. Trigger systems
select some events and ignore others. Reconstruction algorithms turn detector
responses into particle-like objects. Calibration determines whether an
energy scale or momentum measurement can be trusted. The model's public
success is carried by this chain of translation.

That chain also means that a particle is not a bead seen by the eye. A muon
is a track and detector response under a theory of matter and instrumentation.
A jet is a reconstructed expression of quarks and gluons that are not isolated
as free particles. A neutrino often appears as missing momentum. A Higgs
candidate appears through decay products and invariant-mass patterns. The
compiled model teaches the detector what kinds of absences and presences can
count.

The compilation therefore includes null results as real members. A forbidden
decay that remains unseen, a cross section that stays within limits, a search
channel that excludes masses or couplings, and a rare process constrained by
background all shape the model's domain. The public often notices discovery
more than exclusion, but a model is built from both. Absence under a strong
search is a record.

Precision matters because a spectacular event may draw attention, but a
precision mismatch can be more revealing. The anomalous magnetic moment,
electroweak fits, rare decays, and scattering measurements all teach the
community where the compiled structure is tight and where stress may appear.
The model is tested not only by discovering particles but by asking whether
the ordinary numbers stay ordinary at higher precision.

The Higgs boson discovery sits near the end of this compilation not because it
solved everything, but because it completed a crucial mechanism in the working
structure. The 2012 announcements by ATLAS and CMS belong in the discovery
threshold section below, yet their significance belongs here too. The new
particle was not just another bump. It was a record entering a preexisting
compiled structure that had been waiting for that role to be measured.

This waiting is a form of disciplined inference. The theory did not license
unlimited confidence before measurement; it licensed a search. The record
then had to show an excess, channels, masses, rates, and compatibility. When
the excess appeared in independent detectors and channels, the compilation
absorbed it. The model became more complete, but still not omniscient.

The Standard Model also keeps the word \lean{LOGICAL} honest. Its logic is
not merely syllogistic. It is mathematical, statistical, instrumental, and
historical. Gauge symmetries constrain interactions. Conservation laws
restrict events. Detector acceptances condition what can be seen. Statistical
procedures decide whether an excess can be called a particle. Historical
records decide which paths are still open. The logic is distributed across
the compiled record.

The section should therefore avoid two simplifications. It should not say the
model is "just a model," as if compilation were a cosmetic exercise. It should
not say the model is "the truth," as if successful compilation erased the
unknown. A realistic model is a structure that has earned authority by
surviving many records and has earned humility by naming the records it does
not yet contain.

The phrase "earned authority" is doing real work. Authority here does not mean
social prestige. It means the model is difficult to ignore because so many
records have been made to answer to it. Electron-positron collisions, hadron
collisions, neutrino scattering, weak decays, electroweak precision tests,
and QCD phenomena do not all flatter the same weak spot. Their convergence
gives the model weight. Their residual tensions keep it open.

The model's limitations also enter through its very success. A theory precise
enough to predict backgrounds becomes a tool for searching beyond itself. If
an event rate exceeds the Standard Model expectation, one must first ask
whether the detector, background, or analysis is wrong. If the excess
survives, the model's success makes the deviation meaningful. A weak model
cannot produce strong anomalies. A strong model can.

This is one of Decision's most important lessons. A compiled structure gives
communities permission to treat some questions as settled for practical work,
while making unsettled questions sharper. Particle physicists do not reargue
every conservation law before each analysis. They use the compiled structure
to build searches. That use is justified only because the structure remains
answerable to records that could expose its limits.

This is one reason the Standard Model belongs before cosmology. In particle
physics, many compiled records come from controlled or semi-controlled
experimental programs. In cosmology, the object of inference is not
controllable in the same way. The Standard Model teaches compilation under
accelerators and detectors; cosmology teaches compilation under sky, time,
distance, and model extrapolation.

The comparison should not make particle physics look easy. Accelerators give
controlled beams, but the resulting events are statistical, filtered, and
theory-laden. Backgrounds have to be modeled. Rare processes have to be
separated from common ones. Detector upgrades change acceptance. Pileup,
calibration, luminosity, and selection cuts all matter. Control is not the
same as simplicity.

Nor should the comparison make cosmology look weak. A field can be less
controllable and still highly disciplined if its independent records constrain
one another. The point of putting the sections side by side is to show
different forms of compilation. Particle physics compiles under experimental
production. Cosmology compiles under observational reach. Decision asks what
each form licenses.

The Standard Model's history also shows how instrumentation and theory
co-evolve. Better accelerators do not merely provide more of the same events.
They open new energy regimes, new backgrounds, new systematics, and new
possible failures. Better detectors do not simply see more clearly; they
change which signatures can be reconstructed and which questions can be
asked. A compiled model grows by having its theoretical structure repeatedly
translated into new instrumental languages.

Renormalization and gauge structure matter here not because every reader must
follow the mathematics, but because they show that the compilation has formal
discipline. The model is not a curve-fitting warehouse. It is constrained by
mathematical consistency, symmetry, and calculability as well as by measured
events. If the theory could accommodate anything, the compilation would be
weak. It is strong because it cannot.

At the same time, formal discipline does not excuse empirical failure. A
beautiful gauge structure must still meet cross sections, branching ratios,
scattering data, and detector events. The section therefore prepares the
Noether--Galois discussion later. Formal structure can license inference only
when its domain is tied to the measured case. The Standard Model lives by
that tie.

The model's public durability also depends on teaching. Generations of
physicists learn a compressed diagram of particles and interactions. That
diagram is useful, but it can hide the evidential ecology beneath it. The
argument uncompresses the diagram enough to show why it deserves trust. The trust
belongs not to the chart but to the records that made the chart stable.

The strongest compiled structures become boring in use. A background
calculation in an analysis may rely on Standard Model processes without
dramatic explanation. That routine reliance is not laziness when the domain
is appropriate. It is one sign that the compilation has become infrastructure.
The danger is forgetting the domain and treating routine as universal
permission.

This is where beyond-Standard-Model searches get their dignity. They are not
mere dissatisfaction with a successful model. They are attempts to press the
model at its named edges: neutrino mass, dark matter, matter-antimatter
asymmetry, hierarchy questions, gravity, flavor puzzles, and precision
anomalies. A strong model teaches experimenters where to look for its own
limits. Compilation becomes a guide to refusal.

The model also teaches a discipline of scale. Some questions are low-energy
and precision-heavy. Others require higher collision energies, rare-event
searches, astrophysical records, or underground detectors. A compiled model
does not give one experimental route for every question. It organizes a map of
routes. Decision under the model includes deciding which route can actually
bear on the question asked.

This route-choice problem is easy to miss in simplified accounts. If the
question concerns a rare decay, luminosity and background rejection may matter
more than headline energy. If it concerns a heavy new particle, energy reach
may dominate. If it concerns a tiny deviation in a known process, calibration
and theoretical uncertainty may decide the case. The compiled structure does
not eliminate experimental judgment. It makes judgment more exact.

The Standard Model's relation to failure is also mature. Most searches find
nothing beyond expectation. A null result can be disappointing, but it is not
empty. It excludes parameter space, tests backgrounds, improves methods, and
sharpens future design. A field that can learn from null results has moved
beyond spectacle. It has learned decision under absence.

There is a public lesson here for all science. The strongest evidence systems
do not only announce successes. They also publish limits, non-detections,
uncertainties, and exclusions. Those records prevent hope from rewriting the
world. The Standard Model's compiled authority depends as much on disciplined
absence as on famous particles.

The model is therefore an anti-heroic object. It is made by many experiments,
many collaborations, many machines, many failed searches, many calibration
papers, many theoretical repairs, and many ordinary calculations. The named
authors in citations are entry points into that history, not owners of it.
Decision becomes public through the distributed record.

Precision electroweak work shows the compiled character especially well. One
does not merely look for new particles. One measures known relations so
carefully that unknown contributions would disturb them. A small discrepancy
can become a signpost only because the ordinary expectation is known with
discipline. The model's strength is therefore partly the strength of its
ordinary predictions.

QCD gives another kind of compilation. Quarks and gluons are not seen as
free particles in detectors, yet jets, scaling violations, hadronization
models, and high-energy scattering make the strong interaction experimentally
public. The inference depends on theory, detector response, and statistical
patterns. It is a perfect example of realism without naive visibility.

Electroweak unification gives a different lesson. Processes that look unlike
one another at low energy become related under a formal structure. The
records do not become identical; they become mutually interpretable. This is
what compilation does at its best. It lets difference become evidence for a
shared rule without flattening the differences.

The Higgs field then enters not as a decorative final piece but as a mechanism
that makes the compiled structure viable. The discovery of a Higgs-like boson
was a decision under an immense prehistory: predicted role, detector search,
statistical excess, independent channels, property measurements, and later
tests. The public announcement was short because the route was long.

The Standard Model also clarifies how scientific realism can be nonliteral.
No one has to imagine gauge bosons, fields, or quarks as tiny classical
objects in order to take the model seriously. The seriousness comes from
their role in a structure that predicts, constrains, and survives records.
Decision under modern physics often requires this nonliteral realism:
believe the inferential structure where it has earned belief, not a childish
picture of the structure.

The nonliteral realism matters whenever public arguments demand a simple
image before granting belief. Fields, virtual processes, renormalization,
color charge, and symmetry breaking do not become unreal because they resist
cartooning. They become scientifically serious by doing work inside a
compiled structure that repeatedly meets measurement. Decision should not
lower its standard to whatever can be drawn like a little machine.

At the same time, the section refuses the opposite prestige move: hiding
behind mathematics when the record is unsettled. A formal structure that has
not met the relevant measurements does not have the same decision status as a
compiled model. The Standard Model's authority is not mathematical beauty
alone. It is beauty under record pressure.

This double refusal is what makes the model useful for the whole chapter. It
shows how an inference can be abstract, nonliteral, and still deeply real,
while remaining domain-bound and revisable. That is exactly the kind of
realism Decision needs.

\begin{phenom}{The Standard Model Compilation Effect~\cite{glashow1961,higgs1964,weinberg1967,gross1973,politzer1973}}
\label{ph:standard-model-compilation}

\PhStatement
A scientific model becomes compiled when many distinct records constrain one
working structure without collapsing their sources.

\PhOrigin
Electroweak theory, symmetry breaking, quantum chromodynamics, accelerator
experiments, and detector records were assembled across decades into the
modern Standard Model.

\PhObservation
The evidence is heterogeneous: cross sections, decay channels, conservation
rules, particle discoveries, null searches, and precision residues.

\PhConstraint
The compilation may not erase the different strengths, origins, and limits of
the records that support it.

\PhConsequence
The model is inferentially powerful because it is compiled, not because any
single experiment owns the theory.
\end{phenom}

The compilation effect also clarifies what "standard" should mean. Standard
does not mean beyond challenge. It means stable enough to organize work,
precise enough to punish deviations, and incomplete enough to remain a
research program. A standard model becomes useful by giving experimenters
something strong to test against. The strength of the standard is measured
partly by the quality of the failures it can reveal.

That is the Decision lesson. A compiled model can license action inside a
domain: design an experiment, calculate a background, identify a decay,
reject a proposed event, estimate a rate. But every such action borrows the
model's domain. Outside that domain the same authority must weaken. Decision
is the art of knowing where compiled strength applies.

The Standard Model's public role is therefore almost infrastructural. It is a
map on which new experiments are planned, backgrounds are estimated, detector
signatures are classified, and deviations are named. A map can be accurate
for its scale and still omit terrain. The model's incompleteness is not a
reason to discard the map; it is a reason to read the legend.

The compiled model also teaches patience with theory. Some theoretical
structures are not immediately measured. They guide searches, organize
relations, and wait for apparatus to catch up. That waiting is legitimate
only when it remains exposed to measurement. The Higgs mechanism waited in
that disciplined way. Speculation that cannot state how it will meet a record
does not have the same status.

In this section, then, \lean{INFERRED} means something stronger than guess and
weaker than possession. The Standard Model is inferred as a working structure
from a dense ecology of records. Its power comes from density. Its humility
comes from ecology: change the environment, and the model must answer again.

The section belongs in a historical volume rather than a summary of
contemporary physics because the model's authority is historical in the
good sense. It has survived routes through many machines, many analyses, many
failed alternatives, and many ordinary uses. A purely logical summary would
miss that earned weight. A purely anecdotal history would miss the formal
constraint. Decision needs both.

The compiled model finally teaches the reader how to respect a strong
scientific object without worshipping it. Use it where it has earned use.
Test it where it names stress. Do not throw it away because it is incomplete.
Do not call it complete because it is strong. That balance is the decision
posture Decision wants.

In practice, that posture looks ordinary: calculate the expected background,
quote the uncertainty, compare the excess, publish the limit, improve the
detector, and ask again. The grandeur of the model lives inside that
ordinary discipline. Decision is often less theatrical than discovery stories
make it seem.

The model also teaches how a compilation can become a shared language. A
physicist in one experiment can speak of channels, backgrounds, couplings,
decay modes, and exclusions in ways another experiment can understand. This
shared language is not mere culture. It is a transport medium for inference.
Without it, each detector would have local results but no common structure in
which to compare them.

Shared language can become dangerous if it hardens into unexamined habit.
The same names that make work possible can hide assumptions. Background,
signal, efficiency, acceptance, discovery, exclusion, coupling: each word has
a technical route. Decision requires remembering the route when the word is
used to carry public force.

The Standard Model section sits at the center of the scientific arc because it
shows a mature compiled structure being used daily, tested
continually, trusted within scope, and pressed at its edges. That is the form
of inference many sciences envy: powerful enough to guide work, public enough
to be attacked, incomplete enough to keep inquiry alive.

The envy is deserved but should be precise. The model's authority did not
come from one victorious argument. It came from repeated contact between
formal structure and experimental record until many alternatives became too
expensive to keep. A compiled model is what remains after history has made
some possibilities costly.

This cost is intellectual, not merely social. An alternative must explain the
records the Standard Model already explains, recover its successful limits,
and then earn its own additional claim. That is a high price, and it should
be. Strong compilations make rivals do real work. They do not forbid rivals;
they make rivalry answerable.

The same point explains why anomalies are precious. A genuine anomaly is not
just an embarrassment to the compiled structure. It is one of the few places
where a rival may begin to pay the price. Most hints fade; some become
calibrations; a few become openings. The model's strength is what makes the
opening meaningful.

The compiled model therefore makes curiosity more disciplined. It tells the
field which surprises are cheap and which surprises would be expensive enough
to matter. A bump in a poorly understood background is one kind of surprise.
A persistent deviation across calibrated routes is another. Decision depends
on knowing the difference.

The same discipline protects public imagination. It is good for a culture to
imagine physics beyond the Standard Model. It is bad for a culture to treat
every imagined extension as if it had the same status as the compiled record.
Speculation is healthiest when it knows what it owes measurement.

\section{Cosmological Inference}
\episodecoord{1900s onward | Decision :: INFERRED}

The universe as a whole is never in the laboratory. Cosmology is therefore the
place where inference has to be most honest. Expansion records, light-element
abundances, background radiation, supernovae, structure formation, and model
assumptions all carry different weights.

This is the place where bad rhetoric is most tempting. People want cosmology
to deliver a picture of the whole. It does deliver a powerful, public,
mathematical history of the observable universe. But that history is not an
object placed on a bench. It is an inference compiled from light, spectra,
distances, backgrounds, abundances, maps, and relativistic models. The
largeness of the subject makes the route more important, not less.

Expansion evidence begins with a local act of measurement: spectra, redshifts,
distance indicators, calibration, and comparison. The Hubble relation did not
put expansion into a jar. It arranged galaxy distances and velocities into a
pattern that made a static picture less tenable. Later records refined,
challenged, recalibrated, and extended that pattern. The inference grew by
turning the sky into a measured relation.

The distance side of that relation is itself a ladder of trust. Parallax,
standard candles, period-luminosity relations, supernova calibration, and
other methods do not all have the same status. Each rung transfers scale from
nearer and better-known systems to farther and less direct ones. A ladder can
be strong and still require inspection at every rung. Cosmological inference
does not become weaker by admitting the ladder. It becomes realistic.

Redshift is also not a magical distance label. It is a spectral displacement
interpreted through physics and cosmological model. Peculiar velocities,
calibration, selection effects, and the expansion model all enter. The
redshift-distance relation becomes powerful because many such local measures
assemble into a large-scale pattern. The pattern licenses expansion; it does
not let the reader forget spectroscopy.

Light-element abundance gives a different kind of record. It does not look
like expansion on a graph. It asks whether early-universe conditions can
produce the observed proportions of light nuclei. The claim depends on
nuclear physics, observational abundance estimates, astrophysical processing,
and cosmological model assumptions. Its strength is not identical to the
redshift record's strength. The compilation becomes stronger because the
records are different, not because they are interchangeable.

The cosmic microwave background adds transported residue. It is a sky-wide
record with local instrument histories, foreground
subtractions, spectral measurements, anisotropy maps, and parameter fits. It
does not contain the universe by itself. It constrains a universe story when
joined to expansion, nucleosynthesis, structure, and gravitational theory.
Decision asks what that joining licenses.

The CMB is especially strong because it is both simple and structured. A
nearly thermal background gives one kind of evidence; tiny anisotropies give
another; polarization and correlations give further constraints. The same sky
record can support broad history and precision parameter inference. But each
use has its own route. A map used for a public image is not the same as a
likelihood used in parameter estimation. Decision must know which use is
being made.

Foreground subtraction is a good example of strength through difficulty.
Galactic dust, synchrotron emission, point sources, instrument noise, and
scan patterns do not make the CMB unusable. They make the analysis serious.
A cleaned map is not less scientific than a raw map simply because it has
been processed. It is more scientific if the processing route is public and
the uncertainty follows the cleaned signal into inference.

Supernova evidence for acceleration provides another lesson in compiled
inference. Distant standardizable candles, light curves, dust, selection
effects, calibration, and statistical comparison all enter before one speaks
about accelerated expansion or dark energy. The conclusion is powerful
because independent teams and later records made the pattern difficult to
explain away. It remains inferential because every rung has assumptions.

Acceleration also shows how surprise enters a compiled field. The record did
not merely confirm what everyone already wanted. It forced a revision in the
energy content assigned to the cosmos, or at least in the model used to
describe large-scale expansion. The conclusion became acceptable because the
data entered existing and new routes of check: independent supernova teams,
CMB consistency, large-scale structure, and later probes. Surprise had to be
compiled before it became cosmology.

Dark energy language should therefore remain careful. The acceleration record
licenses a term in the model and a problem for physics. It does not hand over
a simple substance. Calling it dark energy is a way to name the inference and
the ignorance together. Decision under cosmology often has this form: a
parameter is measured more securely than its underlying interpretation is
understood.

The word "assumption" should not be used as a sneer. All inference has
assumptions. The honest question is whether the assumptions are public,
testable, constrained by other records, and named at the point where they
enter. Cosmology becomes irresponsible only when assumptions are hidden or
when the conclusion is made to sound stronger than the chain that carries it.

Cosmological inference is a better phrase than cosmic possession. The
observable universe becomes publicly inferable through a disciplined
chain. The chain can be remarkably strong. It can also have uneven thickness.
Direct spectra are not the same kind of record as model-dependent parameters.
Anisotropy maps are not the same kind of record as extrapolations beyond the
observable horizon. Decision's realism depends on keeping these differences
visible.

The same realism protects the argument from both naive skepticism and naive awe.
Naive skepticism says: because cosmology is inferential, it is mere story.
Naive awe says: because the story is magnificent, it has become final truth.
The scientific answer is harder. Some cosmological inferences are among the
best-supported claims humans have about the large-scale universe. They are
best-supported because their routes can be named, checked, corrected, and
combined.

The universe story is therefore a compiled decision under records, not a
simulation verdict. If one asks whether the entire universe is a simulation,
cosmology can constrain some physical stories about what is observed, but it
does not compute the metaphysical proposition from inside the record. What it
can do is say what expansion, background radiation, light elements, structure,
and gravity jointly license about the history of the observable universe.
That is already an immense conclusion.

The word "observable" carries real work. It is not a timid apology. It is the
domain in which the record has routes. To infer beyond that domain requires
additional model commitments. Those commitments may be reasonable, elegant, or
fruitful, but their status must remain named. Decision is not against
extrapolation. It is against smuggling extrapolation into measurement.

Cosmology also shows how a compiled inference can be revised without
collapsing. New distance measurements, tensions in inferred parameters, better
foreground modeling, improved abundance estimates, or alternate probes can
stress the structure. Stress is not the same as failure. A living compiled
field contains tensions because its records have become precise enough to
disagree meaningfully. Decision under cosmology must know whether a tension
is local calibration, new physics, model inadequacy, or ordinary uncertainty.

The current style of cosmological tension is a useful training in humility.
Different routes to a parameter may disagree beyond what their stated
uncertainties comfortably allow. That disagreement may be systematic error,
underestimated uncertainty, or a clue. A responsible field does not instantly
declare revolution, and it does not bury the tension to protect a model. It
keeps the routes visible and asks what each one carried.

This is the same moral as failed lab transport, scaled to the sky. If a value
does not travel cleanly across methods, the failure becomes information. The
inference may shrink, split, or demand new measurement. The universe has not
betrayed the model; the record has discovered that the model or one route
needs more work.

The great strength of modern cosmology is that the records are not all one
kind of thing. Expansion, nucleosynthesis, background radiation, lensing,
supernovae, baryon acoustic patterns, and structure formation do not share the
same instruments or same vulnerabilities. Agreement among heterogeneous
records gives inference its force. Disagreement among them gives inquiry its
next work.

This heterogeneity also limits public simplification. A popular diagram of
cosmic history can be useful, but it hides how many different records have
been placed on one timeline. The diagram is a compiled artifact. The reader
should be able to ask which parts are direct measurement, which are inferred
parameters, which are model-dependent reconstructions, and which are
schematic teaching devices.

The cosmological section therefore ends in a disciplined grandeur. The
universe is not owned. It is inferable. The grandeur is not reduced by that
word. It is made trustworthy by it.

The inference also depends on a careful use of time. Cosmology speaks about
early epochs, recombination, nucleosynthesis, galaxy formation, acceleration,
and future trajectories. These are not memories in the ordinary human sense.
They are ordered reconstructions from records now available. A distant
galaxy's light is old; a background photon carries a different epoch; an
abundance record carries another. The universe story is a time-ordered
compilation, not a diary written by the cosmos.

This time ordering is one reason the "simulation" reduction is too fast. A
simulated world and a non-simulated world could both contain internal records.
The experimental question is what the internal records license. Cosmology can
tell us a disciplined story about the observable universe's history under
physical models. It cannot step outside the record to certify the metaphysical
substrate by declaration. The restraint is what keeps its claims
real.

That restraint also matters for talk about beginnings. Cosmology can describe
hot dense early states, expansion histories, background radiation, and the
limits of current models. It can ask what a singularity in a model means, what
inflation would explain, what quantum gravity might require, and what traces
would count for or against proposed extensions. But the farther the inference
moves from direct record, the more loudly the model commitments must be named.
The beginning of a model is not automatically the beginning of everything.

The same caution applies to multiverse language, cosmic topology, and
speculation beyond the observable horizon. Such ideas may be mathematically
or physically motivated. They may even become indirectly constrained in some
forms. But their decision status differs from expansion, nucleosynthesis, CMB
anisotropy, or supernova acceleration. A realistic cosmology section should
not flatten these into one authority. It should let each claim carry its
route.

This sorting is not meant to make cosmology smaller. It makes its strongest
claims shine more clearly. The expansion of the universe, the relic
background, the broad thermal history, and the evidence for accelerated
expansion are not weakened by refusing to overstate more speculative edges.
They are strengthened because the reader can see which rungs are load-bearing.

Cosmological parameters are another kind of decision object. A number for the
Hubble constant, matter density, baryon fraction, curvature, or amplitude of
fluctuations is not a raw fact in isolation. It is an inference from data and
model. Parameter tables are useful because they compress complicated routes
into comparable values. They are dangerous when the compression hides which
data and assumptions produced them.

The public often encounters cosmology through images: deep fields, background
maps, expansion diagrams, artist renderings. These images are powerful
transport devices. They also risk confusing aesthetic authority with
inference. A map of the microwave sky is a data product with a route. A
timeline is a model summary. A rendering is pedagogy. Decision asks the
reader to admire such images while still asking what record they carry.

The cosmological section therefore contains a lesson about humility under
scale. The larger the subject, the easier it is to smuggle awe into proof.
The argument wants awe, but awe disciplined by route. The observable
universe is not less magnificent because we know it through calibration,
uncertainty, and model comparison. It is more magnificent because such finite
tools can reach so far.

The route also includes selection. Telescopes see what their surveys are able
to see. Flux limits, sky coverage, dust, cadence, wavelength, detector
sensitivity, and analysis pipelines all shape the accessible universe. A
survey record is not a transparent inventory of everything. It is a sampled,
selected, calibrated path through the sky. Cosmological inference has to know
the selection function before it can turn catalogues into conclusions.

Selection is not a flaw if it is public. A well-characterized survey can
support strong inference because the missingness is modeled or bounded. A
poorly characterized survey can make a beautiful catalogue dangerous. This is
the astronomical version of missing data in a trial and detector acceptance in
particle physics. Decision keeps meeting the same form under different names.

The universe story also depends on consistency across epochs. A model that
fits the early background but fails late-time structure is under pressure. A
distance scale that fits local supernovae but strains against background
inference is under pressure. A light-element prediction that misses observed
abundances is under pressure. The strength of cosmology comes from forcing
one story to answer records made at different times and by different means.

This cross-epoch discipline is one reason cosmology can speak so boldly
without touching the whole universe at once. It does not need one impossible
instrument. It uses many finite instruments whose records constrain one
another across time. The inference is large because the comparison is large.
The records remain finite because all science remains finite.

The result is a public cosmos that is neither myth nor possession. It is a
working history under constraints. The history can be taught, tested, revised,
and used. It can inspire without escaping its apparatus. That is the precise
tone this chapter needs before it turns to climate, where public consequence
makes the same inferential discipline more politically volatile.

Cosmology's public consequence is not only philosophical. It changes how
people understand time, matter, origin, scale, and human smallness. That makes
the inference culturally powerful, and cultural power can tempt overstatement.
A responsible cosmology therefore has to be especially clear about what is
measured, what is inferred, what is model-bound, and what is speculative.

The section also has to honor the instruments without making them idols.
Telescopes, spectrographs, satellites, detectors, surveys, and archives are
the routes by which the cosmos becomes public. But no single instrument owns
the cosmic story. Each contributes a window with a wavelength, sensitivity,
resolution, noise profile, and selection. The compilation is stronger than
any window because the windows can be compared.

There is a lesson here about humility before disagreement. When two
cosmological routes disagree, it is not automatically a crisis for science or
a triumph for skeptics. It is a demand for better accounting. Which
calibration moved? Which model assumption mattered? Which uncertainty was
underestimated? Which new physics, if any, would reconcile the records?
Disagreement becomes productive when it is routed.

This routed disagreement is part of why the universe can be publicly
inferable. A private myth cannot be pressured by a parameter tension. A public
cosmology can. It can be made uncomfortable by its own records. That
discomfort is one of its credentials.

The same point holds for model comparison. Competing cosmological models are
not judged by elegance alone. They are judged by what they explain, what they
predict, how many parameters they require, where they fail, and how they fare
against independent records. A model that fits one dataset by becoming
unconstrained everywhere else has not earned the same decision status as a
model that survives multiple routes.

This does not mean simplicity always wins. The universe is not obliged to be
simple. But extra complexity has to answer to records. It must explain
something, predict something, resolve a tension, or open a testable route.
Cosmology is full of beautiful possibilities; Decision asks which
possibilities the compiled record can responsibly carry.

The observable-universe boundary also makes humility concrete. There may be
regions, epochs, or structures beyond present observation. A model may extend
toward them. But extension is not the same as measured fact. Decision
therefore lets cosmology be bold within the observable chain and careful at
the edges where the route thins.

\begin{phenom}{The Cosmological Inference Effect~\cite{alpher1948,hubble1929,penzias1965,perlmutter1999,riess1998}}
\label{ph:cosmological-inference}

\PhStatement
A universe story is licensed only by a compiled chain whose rungs remain
visible.

\PhOrigin
Modern cosmology draws on expansion evidence, primordial nucleosynthesis,
cosmic microwave background observations, distance ladders, supernovae, and
relativistic modeling.

\PhObservation
No single record contains the universe. The conclusion is carried by agreement
among records with different instruments, epochs, and assumptions.

\PhConstraint
Cosmological inference must preserve the difference between direct observation,
calibrated proxy, theoretical bridge, and model-dependent extrapolation.

\PhConsequence
The universe is publicly inferable. It is not possessed by the inference.
\end{phenom}

The phenom's final sentence is the hinge of the argument. Public
inferability is enough to matter. It is enough to organize astronomy, physics,
cosmic history, and future observation. It is not enough to abolish the
difference between record and universe. The distinction is the book's
realism.

Decision can admire cosmology without letting cosmology become a theology of
instruments. The records are strong because they are finite and
public. They carry dates, calibrations, wavelengths, maps, distances,
uncertainties, and models. Their finitude is not the enemy of cosmic thought.
It is the means by which cosmic thought becomes scientific.

The section also prepares the climate section by clarifying model dignity.
Models are neither oracles nor decorations. They are disciplined structures
that let records answer questions records could not answer alone. A model
becomes dangerous when its assumptions vanish from view. It becomes powerful
when assumptions remain visible and the model's consequences meet independent
records.

Cosmology therefore gives Decision its largest honest sentence: the universe
as a whole is not experimentally possessed, but the observable universe is
publicly inferable through a compiled chain. That sentence is modest enough
to be true and large enough to be worth a life of science.

The chain remains open. Better distance measurements, gravitational-wave
standard sirens, improved CMB polarization, surveys of large-scale structure,
and new particle or gravitational physics may change the compilation. That
openness does not make the present inference empty. It makes the inference a
living public object. A living object can be trusted within scope and revised
when the scope changes.

The cosmological decision is therefore exemplary: conclude strongly where the
compiled chain is strong, mark tension where routes disagree, label
speculation where the route thins, and refuse to convert grandeur into
possession. This is the book's cosmology.

That posture makes cosmology more, not less, astonishing. The book does not
need a metaphysical overclaim to honor the achievement. Finite observers have
made an ancient, expanding, structured universe publicly inferable. That
sentence is already large enough.

Cosmology also teaches patience with layered certainty. The statement that
the universe is expanding is not at the same inferential distance as a
particular inflationary model. The statement that the microwave background is
real is not at the same distance as every parameter inferred from a chosen
dataset combination. The statement that acceleration is observed is not the
same as knowing the underlying nature of dark energy. Mature public science
can say all these things at once.

Layered certainty is harder to communicate than a single cosmic slogan, but
it is more truthful. It lets the strongest records carry real authority while
allowing thinner claims to remain provisional. It also lets future records
improve the thin places without making the strong places appear dishonest.

That is why the section keeps returning to the chain. A chain is not only as
weak as its weakest rhetorical flourish. It has rungs with different tasks.
Some support the main inference; others extend, refine, or speculate. The
reader's job is not to admire the chain as one shiny object, but to see which
rungs carry which load.

The load-bearing image is deliberately practical. Cosmology may be sublime,
but its public force comes from knowing which records carry which claims. A
beautiful universe story that cannot name its load-bearing records is not yet
Decision. It is atmosphere.

The same image protects education. Introductory cosmology must simplify, but
it should simplify in a way that leaves a path back to the records. A student
can first learn the timeline, then learn which measurements support each
region of the timeline. The simplified story should be a map, not a mask.

The section's restraint is pedagogical as well as philosophical. It teaches
readers how to admire a cosmic claim without losing the habit of
asking where the claim entered the record. That habit is the difference
between scientific awe and decorative awe.

It also keeps cosmology open to future instruments without humiliating the
present. A future record may refine or revise a claim; that possibility is
part of the claim's public life, not a reason to withhold all conclusion now.
The cosmos is inferred under present routes and left answerable to later
ones.

This is the right posture toward every large cosmological question. Ask what
the present route carries, ask where it thins, and ask what future record
would strengthen or change it. The question is not whether humans have the
whole universe in hand. The question is which parts of the universe have been
made answerable.

That question is humbler and more productive. It turns metaphysical hunger
into experimental work: better distances, cleaner backgrounds, wider surveys,
new messengers, sharper models, and more honest uncertainty.

\section{Discovery Thresholds}
\episodecoord{1900s onward | Decision :: LOGICAL -> HALTED}

A discovery threshold is a community's way of saying where the argument halts
for public use. It does not mean doubt has vanished. It means the record has
crossed a named boundary.

The halt is necessary because inquiry otherwise has no public shape. If every
excess is announced as discovery, the record becomes noise. If no excess can
ever be named until metaphysical certainty arrives, discovery becomes
impossible. A threshold gives the community a rule for moving from search to
announcement under uncertainty. It is a social convention, but not a casual
one. It is tied to error rates, backgrounds, trials factors, detector
channels, and the cost of false discovery.

Five sigma is famous because particle physics learned, painfully, that
interesting bumps are abundant. Background fluctuations, selection choices,
look-elsewhere effects, detector artifacts, and analysis flexibility can all
make the record appear more decisive than it is. A strict threshold is a
public confession of vulnerability. It says: our instruments and analyses can
fool us, so discovery requires a named excess over modeled background.

The threshold does not work alone. It depends on event selection, calibration,
background modeling, control regions, independent channels, collaboration
review, and often independent detectors. ATLAS and CMS did not discover the
Higgs boson by invoking a magic number. They used the number inside a larger
infrastructure of detector understanding and statistical discipline. The
number is the public halt condition for a route.

\lean{HALTED} is one of the book's most easily misunderstood names. It does
not mean stopped forever. It means stopped for this kind of
public conclusion. After the halt, work continues: property measurements,
decay rates, spin-parity tests, coupling fits, searches for deviations, and
new physics beyond the discovered particle. Discovery is a halt in one
argument and an opening in another.

The threshold also protects public language. Before the boundary, one may say
candidate, excess, anomaly, hint, tension, or search result. After the
boundary, one may say discovery under the relevant convention. These words are
not cosmetic. They tell other communities how much weight to place on the
record. Bad public science often begins by stealing the later word too early.

There is no threshold that prevents all error. A five-sigma convention is not
metaphysical armor. It is an engineered public standard under known histories
of error. The standard can be criticized, adjusted, supplemented, or
interpreted differently by field. Its force comes from being named before the
temptation of a particular result, not from being eternal.

The discovery threshold also shows the relation between \lean{LOGICAL} and
\lean{MEASURED}. Logic alone cannot discover a particle. Measurement alone
cannot decide discovery without a rule. The threshold joins them: a measured
excess crosses a logical boundary under a model of background and uncertainty.
The conclusion is licensed because the boundary was public.

This section belongs after cosmology because it gives an explicit form to a
decision that many sciences make implicitly. When is the record enough to
name? When is a pattern a discovery rather than a fluctuation? When does a
community owe the public a conclusion? Particle physics gives one of the
cleanest explicit answers, but the underlying problem is general.

The public word "discovery" also changes incentives. Before announcement, a
collaboration may be cautious internally. After announcement, textbooks,
funding, reputation, and future work reorganize around the named object. A
threshold protects that transition from premature glory. It asks the public
story to wait until the record has crossed a boundary strong enough to bear
the consequences of naming.

Independent detectors make the halt more credible. ATLAS and CMS were not the
same local machine. They shared a collider and a broad theoretical target,
but they had different detector designs, analyses, collaborations, and local
vulnerabilities. Agreement across them made the Higgs conclusion more
transportable. The discovery was not just a bump in one instrument's life.
It was a convergence across partly independent routes.

The threshold also teaches a lesson about background. A signal is meaningful
only against a model of what else could have produced the events. Background
modeling is therefore not a technical appendix to discovery. It is the
negative space that lets the excess have shape. A weak background model can
make a false discovery look strong; a strong background model can make a real
excess publicly visible.

Look-elsewhere effects are another way desire can enter. If a search scans
many masses, channels, cuts, or variables, the chance of some impressive
fluctuation rises. Correcting for that search field is a way to keep the
record from treating an accidental best case as a preordained target. The
community's halt condition must know how many doors it opened.

This is the particle-physics cousin of clinical pre-specification. In both
cases the record becomes dangerous when the question is chosen after the
answer is glimpsed. A named threshold and an analysis route protect the
public word from hindsight. Decision is often the discipline of not letting
afterward pretend to be before.

The threshold also creates a discipline of waiting. Waiting is emotionally
hard in fields built around rare, expensive, spectacular events. A hint may
be real. A hint may also fade. The community's credibility depends on knowing
how to live with the hint until the record can bear the stronger word. This
is not lack of imagination. It is respect for the public force of naming.

Once the threshold is crossed, caution does not disappear. The discovered
object must have properties. Mass, spin, parity, couplings, width, decay
channels, production modes, and consistency with the compiled model all
become new questions. Discovery is therefore a gate, not a throne. It admits
the object into public work.

This gate structure is the reason \lean{HALTED} belongs in the procession.
The halt is local and productive. It stops one uncertainty so other
uncertainties can be pursued more sharply. A community that cannot halt never
builds. A community that halts too easily builds on noise.

The five-sigma convention also shows that thresholds are field-specific. A
clinical trial, bridge rating, climate attribution statement, and particle
discovery do not share one universal number. The standard depends on the
noise field, prior history, consequences, repeatability, and cost of error.
The deeper principle is not five sigma itself. The deeper principle is a
publicly named halt condition appropriate to the decision.

This matters because imported thresholds can become theater. Demanding
particle-physics standards from every field may sound rigorous while making
ordinary decision impossible. Accepting loose standards where false positives
are common may sound flexible while making public language worthless. The
decision craft is to choose and justify the threshold, then live by it before
the tempting result arrives.

Discovery thresholds therefore belong with ethics as well as statistics.
Premature announcement spends public trust. Overcautious silence can delay
knowledge that should guide work. The threshold is a compact between the
community and its future users: this is when we will let a claim carry the
word discovery.

The compact also distributes responsibility. Analysts owe the collaboration a
clean route. Collaborations owe the field a sober announcement. The field
owes the public words that match the record. A threshold is one device by
which these obligations are coordinated. It cannot make everyone honest, but
it gives honesty a public form.

After the threshold, the record still needs transport. Other analyses,
detectors, decay channels, and later runs must carry the object forward. A
discovery that cannot survive property measurement becomes suspicious. A
discovery that survives becomes part of compilation. The halt opens a road.

\begin{phenom}{The Five-Sigma Halt Effect~\cite{atlas2012,cms2012}}
\label{ph:five-sigma-halt}

\PhStatement
A threshold is a public halt condition for inference under uncertainty.

\PhOrigin
Particle physics uses strict statistical standards, including five-sigma
conventions, when announcing discoveries such as the Higgs boson.

\PhObservation
The record includes event selection, background modeling, statistical excess,
look-elsewhere considerations, and independent detector channels.

\PhConstraint
The halt condition must be named before the announcement can borrow the force
of discovery.

\PhConsequence
The threshold licenses public conclusion without pretending that measurement
has become metaphysical certainty.
\end{phenom}

The Higgs announcement is therefore not important only because a particle was
found. It is important because the word "found" was attached to a public halt
condition. That condition let the discovery enter textbooks, talks,
experimental planning, and the Standard Model compilation without pretending
that all future questions had closed.

Thresholds also teach humility about anomalies. A hint below threshold can be
worth studying without being worth announcing as fact. A tension can guide
future work without rewriting the model overnight. A strict threshold protects
both excitement and caution by giving each its proper word.

That vocabulary matters outside physics too. Public cultures are bad at
middle states. They want either proof or nothing. Discovery thresholds model a
healthier middle: promising, suggestive, anomalous, excluded, discovered,
measured. A mature science needs all of these words because the record moves
through stages before decision hardens.

This vocabulary is one of Decision's gifts. It lets excitement remain
honest. One can be thrilled by a hint without calling it a discovery, and one
can be confident in a discovery without calling it final metaphysics. The
words keep the record from being crushed by our impatience.

The halt condition therefore protects time. It gives the community time to
check before naming and gives later work a named place to begin after naming.
It is not a wall. It is a gate with a label.

The label matters because future users need to know what was halted. Discovery
halts the existence question under a threshold; it does not halt property
measurement, interpretation, or future stress.

A discovery paper is a beginning archive. It stores the route by which the name
became admissible, so later work can ask different questions
without reopening everything at once.

The archive function is why methods sections, auxiliary plots, and systematic
uncertainties matter. They are the handle by which the halt can later be
checked.

\section{Climate Compilation}
\episodecoord{1800s onward | Decision :: COMPILED -> INFERRED}

Decision returns to climate only after separating the overhead
problem from the inference problem. The inference section must ask a narrower
question: what does the measured and modeled chain license, and where does the
chain stop?

This section should not repeat Overhead's institutional argument. Overhead
asked how consensus, policy, funding, communication, and propaganda attach to
the climate record. Decision asks what inference the record itself licenses.
Those questions touch, but they are not the same. Confusing them is one of
the main ways public argument becomes dishonest.

The climate record is heterogeneous. Radiative physics gives a mechanism.
Instrumental temperature records give measured trends. Atmospheric
measurements give greenhouse gas concentrations. Ice cores and other archives
extend some variables into the past. Ocean heat content, glacier mass,
phenology, sea level, and satellite records add different lines of evidence.
Detection and attribution studies ask how much of observed change is
consistent with known forcings. Model ensembles project futures under named
scenarios. Each link has its own evidential character.

The strongest public inference is not that every detail is equally direct. It
is that the chain is strong enough to license the basic factual conclusion:
human activity has warmed the climate system, chiefly through greenhouse gas
emissions, and continued emissions alter future risk. That conclusion is not
made weak by admitting uncertainty. It is made scientific by quantifying
uncertainty and naming where the chain is thick or thin.

The chain is thickest where measurement is direct and repeated: atmospheric
CO2 concentration, many instrumental temperature records, satellite radiance
records, ocean heat measurements, and laboratory radiative physics. It is
thinner where the record depends on proxies, reconstructions, sparse
historical coverage, or model-dependent translation. Thin does not mean
worthless. It means the inference must carry more visible uncertainty and
more dependence on method.

This is where public rhetoric often fails. One side may point to uncertainty
in proxy reconstructions or model projections as if uncertainty erased the
measured warming trend. It does not. Another side may treat the measured
warming trend as if it directly dictated every policy instrument, discount
rate, burden-sharing scheme, or technological pathway. It does not. The
record licenses factual inference and risk assessment. Policy choices add
values, costs, tradeoffs, distribution, responsibility, and political
judgment.

The distinction is not a retreat from science. It is respect for science. If
measured facts are used to smuggle in value judgments, the public learns to
distrust the facts. If value disagreements are used to deny measured facts,
the public loses contact with reality. Decision under climate must hold both
boundaries: do not deny what the record licenses; do not pretend the record
alone has chosen the policy.

Scenarios are a useful example. A projection under a high-emissions or
low-emissions pathway is not a prophecy. It is a conditional inference:
under these assumptions about emissions, land use, aerosols, socioeconomic
development, model response, and feedbacks, these ranges follow. The
conditional form is a strength, not an embarrassment. It tells decision makers
which futures are being compared. The dishonesty begins when a scenario is
sold as fate or dismissed as fiction because it is conditional.

Attribution is another example. The question is not merely whether climate
changes naturally. It does. The question is whether the observed pattern,
magnitude, and physics of recent warming are explainable without human
forcing, and how different forcings combine. This is a compiled inference
from physics, observations, and models. It is not a single thermometer
shouting politics. It is a chain of comparison in which alternative accounts
must pay for the record.

The unevenness of the chain should be visible in prose. Direct atmospheric
measurement is not the same as paleoclimate albedo reconstruction. A global
mean temperature anomaly is not the same as local weather damage. A detection
and attribution statement is not the same as a cost-benefit analysis. A model
ensemble range is not the same as an ethical distribution of burdens. Keeping
these differences visible makes the inference stronger, because it prevents
the strong parts from being weakened by overclaim in the weaker parts.

Climate is therefore a severe test of Decision. The record is strong enough
to demand public inference, but the action space is value-laden enough to
forbid measurement from pretending to be sovereign. The correct conclusion is
neither paralysis nor technocracy. It is disciplined public judgment under a
compiled record.

The section also needs to say what denial does. Denial often hides inside a
selective demand for impossible certainty. It says: because projections have
ranges, warming is not real; because proxies require inference, direct
measurements can be ignored; because policy is contested, attribution is
optional. This is not humility. It is a refusal to let measured facts bind
belief. Decision refuses that failure just as firmly as it refuses policy
smuggling.

The climate record therefore lands on Decision's central word:
\lean{INFERRED}. Warming, attribution, and risk are inferred under a compiled
chain whose rungs are public. The inference is strong enough to constrain
responsible action. The choice among actions remains a decision made under
science, economics, ethics, law, politics, and risk tolerance. The record does
not vanish at the policy boundary; it also does not become the whole policy.

Risk is the bridge word. Science can describe hazards, exposure,
vulnerability, likelihoods, confidence ranges, and physical consequences.
Decision then asks how much risk is acceptable, who bears it, who pays to
reduce it, which losses count, and how present costs are weighed against
future harms. Those are not measurements in disguise. They are public
judgments under measured risk.

Mitigation and adaptation show the handoff. Emissions reductions, energy
systems, land use, infrastructure adaptation, insurance, relocation,
resilience, carbon removal, and research into sunlight reflection are not the
same kind of act. They carry different evidence bases, costs, risks,
distributional effects, reversibilities, and governance problems. Climate
science can constrain their physical premises. It cannot make their ethical
tradeoffs disappear.

Albedo is a useful example of chain thickness. Satellite observations can
measure reflected radiation in one way. Paleoclimate albedo reconstruction
requires proxies and model bridges. Proposed reflecting-sunlight interventions
add engineering, governance, termination, regional effects, and moral hazard.
These belong in one broad climate conversation, but they do not have the same
evidential status. Decision becomes honest by refusing to blur them.

The climate section must also protect the word "model." A climate model is not
a crystal ball. It is an organized physical computation constrained by
observations, theory, and comparison. Models can be wrong in detail and still
useful; useful and still uncertain; uncertain and still decision-relevant.
The public mistake is to demand either perfect prophecy or no inference. Real
decision almost always lives between those poles.

Attribution of extreme events gives another boundary. A heat wave, flood,
fire season, drought, or storm is local and weather-shaped. Attribution asks
how the probability or intensity of such events changes under a warmed
climate. That is a conditional statistical inference, not a simple sentence
that climate change "caused" the event in the same way a match causes flame.
The nuance strengthens the record because it names the causal form.

There is also a temporal asymmetry. Some climate facts are already measured:
greenhouse gas concentrations, warming trends, ocean heat increase, ice loss,
sea level rise. Some consequences are emerging. Some are projected under
future emissions and response pathways. Decision must not treat future ranges
as present measurements, but it also must not ignore them because they are
future. Risk decisions are precisely about acting before all consequences
have arrived.

The bridge and treatment vignettes matter after climate because engineers do
not wait for collapse to decide a bridge is unsafe.
Physicians do not wait for every future patient outcome before treating under
trial evidence. Climate decision similarly acts under records that are strong
but finite. The hard question is not whether finite records can matter. They
must. The hard question is how openly the finitude is carried into action.

The public also needs a language for confidence. "Likely," "very likely,"
"high confidence," and quantified ranges are not evasions when used honestly.
They are the way uncertainty becomes usable. Propaganda tends to hate such
language because it wants either total certainty or total doubt. Science needs
graded confidence because decision needs to know how heavily the record
weighs.

This graded language should not be confused with weakness. A bridge engineer
using safety factors is not weak because the future load is uncertain. A
physician using confidence intervals is not weak because the next patient is
not identical to the trial population. A climate scientist using uncertainty
ranges is not weak because the future depends on human choices and physical
feedbacks. Bounds are part of responsible inference.

Climate also forces Decision to name irreversibility. Some decisions can
be revised cheaply. Some cannot. Emitted carbon, lost ice, changed ecosystems,
built infrastructure, and delayed adaptation have temporal inertia. Decision
under uncertainty changes when waiting changes the option set. This does not
answer every policy question, but it shows why uncertainty can increase the
need for decision rather than suspend it.

The climate inference is therefore morally serious without being morally
self-executing. It tells the public that the physical record has weight. It
does not relieve the public of choosing among imperfect responses. That is a
harder message than either denial or technocratic command, and it is the only
message this chapter can defend.

The chapter should also keep local and global claims distinct. A global mean
temperature trend is not the experience of a farmer, coastal city, forest, or
grid operator. Local impacts require regional records, exposure, vulnerability,
adaptation capacity, and social conditions. The global record constrains the
local conversation, but it does not replace local decision. Truth has to
transport again.

Loss and damage language makes this distinction acute. A physical hazard
becomes a human loss through infrastructure, wealth, governance, geography,
history, and vulnerability. Climate science can identify changing hazards and
risks. Public decision must then address responsibility and repair. Those
ethical and political questions should not be hidden inside temperature
graphs, and the graphs should not be denied because the ethical questions are
hard.

The same distinction applies to uncertainty about technology. Carbon removal,
adaptation infrastructure, nuclear power, renewables, demand reduction,
geoengineering research, and other responses have different evidence bases and
failure modes. The physical climate record can make delay risky, but it does
not make every proposed technology wise. Decision needs case-specific records
after the climate inference has done its work.

Economic modeling adds another layer. Discount rates, damage functions,
innovation assumptions, adaptation costs, and distributional weights are not
read directly from thermometers. They are decision tools built on physical
inference and social assumptions. They may be necessary, but their value
choices should be visible. Hiding them under "the science says" harms both
science and public ethics.

Communication is part of decision too. A scientist speaking to the public
must decide how to preserve uncertainty without giving denial a theatrical
opening. This is difficult. Too much compression can become overclaim. Too
much caveat can let bad-faith actors pretend the record says nothing. The
answer is not perfect wording. It is disciplined sorting, repeated patiently.

Climate therefore closes the major scientific sections with a hard lesson:
the more consequential the inference, the more aggressively people will try
to steal or erase its authority. The record must be strong, but it must also
be well-described. Description is defense.

\begin{phenom}{The Climate Compilation Effect~\cite{arrhenius1896,ipcc2021}}
\label{ph:climate-compilation}

\PhStatement
A heterogeneous climate record can license factual inference while leaving
policy choice visibly outside the measurement itself.

\PhOrigin
Climate science compiles radiative physics, instrumental temperature records,
greenhouse gas measurements, paleoclimate proxies, model ensembles, and impact
studies.

\PhObservation
The chain has uneven thickness. Direct atmospheric measurements differ from
proxy reconstructions, and projections under scenarios differ from already
measured trends.

\PhConstraint
Honest inference must state which parts of the chain are measured, which are
proxied, which are modeled, and which are value-laden decisions.

\PhConsequence
Uncertainty does not license denial of measured facts. Measured facts do not by
themselves license every proposed action. Both boundaries must remain public.
\end{phenom}

This phenom is one of the places where the argument must be more realistic
than the opening volume. It must not sound as if the author is choosing a team in the public
shouting match. It must sound as if the record is being sorted. Direct
measurement, proxy reconstruction, attribution, projection, impact, and policy
are related, but they are not identical. The sorting is the argument.

Once sorted, the conclusion can be forceful without being inflated. Human
caused warming is not a vibes claim. It is a compiled inference from physics
and records. Policy disagreement is not automatically science denial. It
becomes denial when it rejects the licensed factual inference. Policy advocacy
is not automatically science. It becomes honest when it states the values and
tradeoffs added after the factual inference. Climate decision needs both
sentences.

The section's final realism is procedural. A public climate decision should
say: here are the measured facts; here are the attribution claims; here are
the projection scenarios; here are the uncertainty ranges; here are the value
judgments; here are the distributional consequences; here is the action
chosen. That order will not remove conflict. It will make the conflict
legible. Legible conflict is better than counterfeit certainty.

The procedure also makes revision possible. If a model range narrows, a
technology improves, a cost changes, an impact arrives sooner, or a policy
fails, the decision can be updated without pretending the entire record has
collapsed. Public trust depends on that distinction between revision and
reversal.

That distinction is especially important because climate decisions unfold
over decades. A policy may be physically sensible and politically poor. A
technology may be promising and socially unacceptable. A projection may be
right in trend and wrong in regional detail. Decision must be able to learn
without letting every update become a referendum on the existence of warming.

Climate compilation thus becomes Decision's sternest public case. The
facts are strong enough to bind; the choices are difficult enough to require
open values; the future is uncertain enough to demand ranges; and delay is
consequential enough to make evasion a decision too.

That last clause is the hard one. In a changing system, refusal to decide is
not neutrality. It is one policy pathway under another name.

The record does not choose the pathway, but it can remove the innocence of
pretending no pathway was chosen.

That is enough to make the inference politically uncomfortable and
scientifically necessary.

\section{Formal Closure and Licensed Inference}
\episodecoord{1800s-1900s | Decision :: LOGICAL seed}

Noether and Galois belong in Decision as formal seeds. They are not
experimental episodes in the ordinary sense, but they show a kind of closure
that experimental inference often borrows. Galois turns solvability into a
structural question about groups. Noether turns symmetry into conservation.
In both cases a conclusion is licensed not by inspecting every instance, but
by finding the structure that governs the instances.

The point is not to collapse mathematics into experiment. It is to show why
\lean{LOGICAL} appears in a book about measurement. Experiment often needs a
formal rule that says what would count as closure, impossibility, symmetry,
or conservation. A detector record becomes meaningful partly because energy,
momentum, charge, spin, or other conserved quantities constrain what can be
inferred. A polynomial problem becomes soluble or insoluble because its
structure permits or forbids a route. Formal closure teaches decision how to
stop without mere exhaustion.

Galois is a useful warning against computation as brute force. One may try
examples forever and still not know the structural condition for solution.
The Galois move asks for the group-theoretic shape that licenses a decision
about solvability. That is a \lean{HALTED} lesson in pure form: the argument
can stop because the structure has spoken, not because the calculator grew
tired.

Noether gives the complementary lesson for physics. A conservation law is not
merely an observed habit. It is tied to symmetry. The formal relation lets
experimental records carry more inference than a list of events could carry
alone. If a proposed event violates a conservation rule without an admissible
account, the record has a reason to reject or reinterpret it. The decision is
logical and measured at once.

Conservation laws feel decisive in experimental practice because they let a
local record say more than "something odd happened." Missing energy
may suggest an unobserved particle only because conservation gives absence a
formal role. A forbidden decay may count as evidence against a proposal only
because the formal rule states what should be conserved. Noether's lesson is
not decorative mathematics. It is one of the ways absence becomes inferential.

Galois gives the same dignity to impossibility. A problem may be impossible
by radicals not because no one has tried hard enough, but because the
structure forbids that route. This matters for the user's concern about what
cannot be computed. Some propositions cannot be settled by running longer.
The decision has to come from understanding the structure of the problem, or
from admitting that the structure does not license the desired conclusion.

Formal closure is therefore a guard against magical persistence. It prevents
the culture from saying, "Perhaps if we compute harder, measure longer, or
desire more intensely, the forbidden route will open." Sometimes more work is
needed. Sometimes the route is formally closed. Decision needs to know the
difference.

The section belongs late in Decision because formal beauty is dangerous if
it arrives too early. A beautiful structure can seduce experiment into
finding what it wants. But after five chapters of marks, residues, staging,
overhead, and transport, formal closure can be given its due. It is one of
the ways public inference becomes more than accumulation.

The due is large but bounded. Formal structure can decide within its domain
with a purity no empirical record can match. But the application of that
structure to a physical case still requires measurement, modeling, and
interpretation. Noether's theorem does not by itself calibrate a detector.
Galois theory does not by itself decide a bridge rating. Logic supplies
closure where the premises and domain are secure. Decision must secure them.

This boundedness is why formal sections are seeds rather than substitutes.
They teach the shape of legitimate stopping: here is the invariant, here is
the symmetry, here is the group, here is the forbidden route. But the world
still has to be shown to instantiate the structure. In physics, that showing
requires apparatus. In public decisions, it requires facts and rules. In
mathematics, it requires proof that the structure has been correctly named.

The formal seed also disciplines rhetoric about impossibility. "We cannot
compute it" is not always the same as "it is not real," and "we assume it is
real" is not the same as "therefore it is true." Galois and Noether teach a
better route: identify the structure that licenses the conclusion. Without
that structure, persistence of assertion is not proof.

This is the formal version of Decision's demand. A decision is not
licensed by intensity. It is licensed by a rule meeting a record or a
structure meeting a case. Formal beauty can help identify the rule; it cannot
replace the meeting.

The formal seed also explains why Decision does not end in computation.
Computation can explore consequences once the structure and inputs are known.
It cannot manufacture the warrant for assuming what had to be proved. Decision
requires warrant before calculation can carry authority.

Formal closure feels severe because it does not flatter desire; it asks
what the structure permits.

That severity is useful in a chapter otherwise full of judgment under
uncertainty. It reminds the reader that some closures are exact where the
domain is exact, and that exactness should be respected without being
exported beyond its domain.

The formal seed is a discipline of scope.

\begin{phenom}{The Noether--Galois Closure Effect~\cite{noether1918,galois1846v6}}
\label{ph:noether-galois-closure}

\PhStatement
Formal structure can license a halt in inference when it states what routes
are possible, impossible, conserved, or forbidden.

\PhOrigin
Galois theory related solvability to group structure, while Noether's theorem
related continuous symmetries to conservation laws.

\PhObservation
The record consists of structural relations: group actions, symmetries,
invariants, conservation rules, and consequences for admissible conclusions.

\PhConstraint
Formal closure must not be mistaken for empirical possession; it licenses
inference only where the structure applies to the measured case.

\PhConsequence
Logic enters experiment as a public rule for stopping, rejecting, or
organizing conclusions under structure.
\end{phenom}

Noether--Galois gives Decision a clean internal skeleton. It says that some
decisions are not made by piling up more cases. They are made by recognizing
the rule that governs the cases. Experimental history needs that skeleton,
but only after the skeleton has been attached to records.

\section{Jury Verdict: Inference Under Instructions}
\episodecoord{ordinary | Decision :: INFERRED}

A jury verdict is not science, but it is a disciplined public inference under
a record. Witnesses, exhibits, expert testimony, instructions, burdens of
proof, and deliberation enter a form that must end in a decision. The legal
system knows something science also knows: local events do not walk into
public judgment by themselves.

The analogy is useful because a verdict is explicitly not the event. It is a
licensed conclusion about the event under admissible evidence and legal
instruction. A jury may be wrong. Rules may exclude evidence. Witnesses may
mislead. Yet the system still needs a public decision. The verdict is a halt
condition under a social record.

This differs from science because legal finality has institutional force.
Science can often reopen more easily. But the parallel clarifies Decision's
central grammar: evidence enters rules; rules license inference; inference
does not become omniscience. The jury shows that public life often acts
without final possession of truth.

The burden of proof is the key decision device. Civil and criminal standards
do not ask for the same strength of record because the consequences differ.
This is a public admission that inference depends on what is at stake. Science
has analogues in discovery thresholds, safety margins, and clinical risk
benefit judgments. The record's weight is judged partly by the cost of being
wrong.

Jury instructions also show why rules must be explicit before conclusion.
Jurors are told what question they are deciding, what counts legally, and what
standard applies. If those instructions are confused, the verdict's authority
weakens. Decision in science has the same need for stated questions. A record
cannot license a conclusion if the conclusion was never the question.

The verdict also shows the pressure of closure. A jury cannot usually publish
"interesting but inconclusive" and send everyone home. It must decide under a
standard. That pressure makes legal inference different from scientific
inquiry, but it also reveals something science sometimes hides: public
institutions often need conclusions before perfect knowledge arrives.

That need is why standards differ. Beyond reasonable doubt is not a law of
nature. It is a public rule fitted to consequence. Decision is always partly
about the cost of error.

The jury vignette keeps Decision social: inference is something
communities do under forms, not only something minds do in private.

The forms may be imperfect, but without them public judgment becomes charisma.

The jury vignette therefore makes rule-bound inference feel civic rather than
merely technical.

It shows that finite evidence and public consequence have always needed
forms.

\begin{phenom}{The Jury-Verdict Inference Effect~\cite{wigmore1904}}
\label{ph:jury-verdict-inference}

\PhStatement
A public verdict is an inference licensed by admissible record and governing
instructions, not possession of the event itself.

\PhOrigin
Legal evidence traditions organize witness, admissibility, burden of proof,
instruction, and deliberation into a form capable of public judgment.

\PhObservation
The record includes testimony, exhibits, expert claims, objections,
instructions, deliberation, and a verdict.

\PhConstraint
The verdict may not pretend to exceed the record and rules that licensed it.

\PhConsequence
Ordinary public judgment shows the shape of inference under finite evidence:
necessary, rule-bound, consequential, and fallible.
\end{phenom}

\section{Bridge Safety Inference: Measurement Under Load}
\episodecoord{ordinary | Decision :: MEASURED -> INFERRED}

A bridge safety decision is a plain example of inference under measured risk.
Inspectors, engineers, load ratings, material records, corrosion, cracks,
fatigue, design assumptions, traffic patterns, and repair histories all enter
the question: may this bridge remain open under these loads? The answer is
not a metaphysical truth about the bridge. It is a licensed public decision
under a record.

The bridge case matters because the decision cannot wait for perfect
knowledge. Closing a bridge has costs. Leaving an unsafe bridge open has
costs. The record must be strong enough to justify action even when hidden
defects, future loads, weather, and aging remain uncertain. Engineering
decision therefore lives in the mature middle of this chapter: measured,
compiled, inferred.

Load rating is not merely arithmetic. It depends on design documents, current
condition, inspection quality, material assumptions, safety factors, legal
standards, and consequences of failure. A crack is a mark; corrosion is a
record; traffic is overhead; inspection transports local witness; the rating
licenses decision. The whole argument is compressed into ordinary
infrastructure.

Safety factors are a public philosophy in numerical form. They say that
uncertainty, variability, deterioration, and consequence require margin. The
margin is not cowardice and not proof. It is a decision rule for finite
knowledge under load. When the margin disappears, the bridge may have to
close even before collapse supplies the final evidence no one wants.

The bridge also shows how records age. An inspection that was adequate five
years ago may not be adequate after floods, heavier traffic, corrosion, or
new cracking. Decision under infrastructure is continuous. The record must be
maintained because the object changes while the public keeps using it.

This makes inspection a form of truth transport through time. The bridge does
not remain the bridge described in the original plans. Weather, load, repair,
and material fatigue rewrite it. A safety decision must transport the design
record into the present condition before it can infer future use.

The bridge vignette therefore gives \lean{MEASURED} a civic face. Measurement
is not only discovery. It is maintenance of public trust under load.

Its modesty is the point: civilization depends on unglamorous inferences made
well.

\begin{phenom}{The Bridge-Safety Inference Effect~\cite{aashto2020}}
\label{ph:bridge-safety-inference}

\PhStatement
Engineering safety decisions infer admissible use from measured condition,
design assumptions, and public load standards.

\PhOrigin
Bridge evaluation practice compiles inspection, material, structural, and load
records to rate bridges for use, repair, restriction, or closure.

\PhObservation
The record includes design, age, materials, inspection, defects, load paths,
traffic assumptions, rating factors, and maintenance history.

\PhConstraint
The public decision must carry both measured condition and safety margin; it
may not treat either uncertainty or convenience as proof.

\PhConsequence
Infrastructure decision shows that finite records can be binding when public
risk requires action.
\end{phenom}

\section{Treatment Decision: Evidence Meets the Patient}
\episodecoord{ordinary | Decision :: LOGICAL -> INFERRED}

A treatment decision is where trial evidence meets a particular person. The
record may show an average benefit in a population. The patient has age,
history, preferences, comorbidities, access, risks, costs, and values. The
clinician's decision is not a rejection of the trial record and not a
mechanical copy of it. It is an inference under evidence, judgment, and care.

Evidence-based medicine is often caricatured as cookbook medicine. Its better
form is the opposite. It asks for the best available external evidence, the
clinician's expertise, and the patient's values to meet openly. That meeting
is Decision in miniature. Measurement enters; logic organizes; uncertainty
remains; action occurs.

The treatment case also returns to the user's hardest proposition: the
uncomputable cannot be made true by assuming it is real. A patient is not a
computable object exhausted by a trial table. But the table still matters.
The decision is neither pure computation nor pure intuition. It is a licensed
inference made under records that do not own the person.

Shared decision is not a sentimental add-on. A treatment may reduce risk on
average while carrying side effects one patient refuses, costs
another cannot bear, or burdens a third cannot practically manage. Values do
not erase evidence; they tell evidence what kind of action it is being asked
to support. The clinician's task is to keep both in view without letting one
pretend to be the other.

Follow-up completes the decision. A treatment inference is tested again in
the person's outcome, adverse effects, adherence, and changing condition. The
decision may be revised. That revision is not failure of evidence-based
practice. It is evidence-based practice continuing after the first action.

The treatment vignette therefore restores the human scale after the large
scientific compilations. The same movement that speaks of particles,
cosmology, and climate ends with a person deciding what to take, refuse,
continue, or stop. This is where inference becomes lived consequence.

It also reminds Decision that uncertainty is not an abstraction to the
person who must act. Waiting, trying, refusing, and stopping all have costs.
The best decision is not the one without uncertainty, but the one that carries
uncertainty honestly into care.

The patient brings the book back to scale. Inference may be compiled by
institutions, but it lands in a body.

That landing is where responsibility stops being abstract.

The treatment decision is small only in scale. In consequence, it can be as
serious as any cosmological conclusion.

The person is the final locality.

\begin{phenom}{The Treatment-Decision Effect~\cite{sackett1996}}
\label{ph:treatment-decision}

\PhStatement
A clinical decision translates population evidence into patient action without
pretending the population record owns the patient.

\PhOrigin
Evidence-based medicine framed clinical practice as the integration of best
research evidence, clinical expertise, and patient values.

\PhObservation
The record includes trial evidence, diagnosis, prognosis, adverse effects,
patient circumstances, preferences, alternatives, and follow-up.

\PhConstraint
The decision must preserve the difference between measured population effect,
individual applicability, and value-laden choice.

\PhConsequence
Treatment shows inference at human scale: strong enough to guide action,
humble enough to remain answerable to the person and outcome.
\end{phenom}

\begin{movementclose}
The procession ends at \lean{INFERRED}. The word is quiet because the record is
doing the work. An inferred conclusion can be strong enough to act on, publish,
build with, treat with, or regulate by. It is not guessed. It is not owned. It
is licensed by a record that remains open to challenge.

This quiet ending should feel earned. Entry taught the book that marks need
custody before they become records. Residue taught it that leftovers can
become disciplined questions instead of embarrassments. Staging taught it that
values are produced by procedures, not plucked from the world. Overhead taught
it that institutions, cost, desire, and authority can carry or corrupt the
record. Truth Transport taught it that local witness must travel with its
conditions intact. Decision has now asked what conclusions such records can
bear.

The answer is neither skepticism nor triumph. The record can bear a great
deal. It can license treatment, bridge closure, particle discovery, climate
attribution, cosmological history, and formal refusal. But the record does not
become the universe, the patient, the city, the future, or the moral law. It
remains a public route to inference. That route is strong because it is
visible.

\lean{HALTED} belongs before \lean{INFERRED}. A halt condition is not a claim
that thinking has ended. It is a point at which a community may
use a conclusion without pretending every possible question has been
answered. The double-blind trial halts one argument about effect. The
five-sigma threshold halts one argument about discovery. A bridge rating
halts one argument about admissible load. Each halt is local to a decision.
Each leaves future records possible.

\lean{COMPILED} is equally important. Modern inference rarely rests on one
record. It rests on many records whose vulnerabilities differ. The Standard
Model compiles accelerator and precision evidence into a working structure.
Cosmology compiles expansion, background radiation, abundances, and
supernovae into a universe story. Climate compiles physics, observations,
proxies, and models into factual attribution and risk. Compilation is powerful
because one record can check another. It is dangerous when the differences
among records are forgotten.

Decision has also shown that ordinary decisions are not lesser
philosophy. A jury verdict, a bridge closure, and a treatment choice make the
same structure visible without the prestige aura of cosmic or particle
science. They decide under finite evidence. They carry consequences. They can
be wrong. They need rules. They remind the reader that inference is a public
practice before it is an abstract noun.

The relation to Volume 5 should now be clear. The names are not decorative
labels. They are obligations. \lean{LOGICAL} asks the decision to state its
rule. \lean{MEASURED} asks it to state its record. \lean{HALTED} asks it to
state the point at which conclusion becomes admissible. \lean{COMPILED} asks
it to preserve heterogeneous sources. \lean{INFERRED} asks it to remain
answerable to the route. A decision that cannot do these things should not
borrow the authority of science.

This is also the final answer to placeless truth. The book has not found a view
from nowhere. It has found something better for finite creatures: routes
from place to public inference. A flask, a plate, a detector, a trial, a
station, a model, a threshold, and a bridge rating are all finite. Their
finitude is not a scandal. It is the condition under which public knowledge
can exist.

The final humility is not weakness. An inferred conclusion can still be
binding. A physician may treat. An engineer may close a bridge. A
collaboration may announce a particle. A climate scientist may state
attribution. A cosmologist may teach a history of the observable universe.
The humility lies in saying what licenses the conclusion and what does not.

That humility is also what keeps the book from becoming propaganda for
science. Propaganda wants the conclusion without the route. Cynicism wants the
route's limits to erase the conclusion. Decision refuses both. It says: here
is the route, here is the licensed conclusion, here are the limits, here is
what remains for future record or public judgment.

This refusal is the book's answer to premature QED. The letters are tempting
because they promise an ending. But experimental history rarely ends by
annihilating all alternatives from outside the system. It ends provisionally,
publicly, and powerfully: enough to treat, build, name, attribute, compile,
or infer; not enough to stop the universe from surprising us. The proof is a
route, not a throne.

Decision also leaves room for courage. Humility should not become an
excuse for evasion. A record that licenses action asks something of us. If a
bridge is unsafe, if a treatment helps, if a particle is discovered, if human
activity warms the climate, if a model reliably organizes a domain, then
refusing all conclusion is not sophistication. It is another way of fleeing
the record.

Likewise, courage should not become swagger. Acting under a record does not
make the actor identical with truth. It creates responsibility to update,
monitor, disclose uncertainty, and distinguish evidence from value. A
decision made under science remains answerable after it is made. The record
does not retire.

That is why the book ends with a public rather than private universe. A
private revelation can be absolute for the person who claims it. A public
inference must carry marks, residue, staging, overhead, transport, and
decision. It is slower, less glamorous, and more humane. It lets other people
check the route.

The checking is not an insult to the conclusion. It is the condition of the
conclusion's dignity. A particle discovery that cannot be checked becomes a
rumor. A treatment result that cannot be checked becomes marketing. A climate
claim that cannot be checked becomes ideology. A cosmology that cannot be
checked becomes myth. Public inference stays alive by remaining vulnerable to
competent challenge.

The final posture is therefore neither defensive nor imperial. It
does not apologize for science by shrinking every conclusion into maybe. It
does not inflate science by pretending every conclusion is final. It says:
the route is good enough here, the route is thin there, the route is missing
over there, and where the route is good enough we may infer.

The phrase "good enough" is not casual. It is calibrated to consequence. Good
enough for a classroom diagram may not be good enough for a drug approval.
Good enough for a search hint may not be good enough for discovery. Good
enough for a bridge inspection last year may not be good enough after a flood.
Decision is the practice of matching evidential strength to the weight placed
on it.

That is the universe here: not the object in itself, and not a
simulation of the object, but the universe as experimental history has made it
publicly inferable.

The last phrase matters. Publicly inferable does not mean finally possessed.
It means open to challenge, repair, repetition where possible, transport where
necessary, and decision where required. The universe is real here in
the only way experiment can give it to us: through records that survive enough
discipline to bear inference.

That is not a consolation prize. It is the central achievement. Finite beings
inside the world have made parts of the world answerable to shared records.
They have done it with flasks, plates, detectors, trials, stations,
equations, thresholds, and ordinary judgments under risk. The procession does
not end by leaving finitude behind. It ends by showing what finitude can do
when it keeps faith with its routes.

So the last word is not certainty, and it is not doubt. It is inference. That
word has work inside it: record, rule, scope, uncertainty, consequence,
revision. It is quiet enough to remain honest and strong enough to carry a
world.

The book can end there because inference is not a lesser truth. It is truth
as finite public beings can responsibly use it. That is enough.

It is also more realistic than a louder ending. A louder ending would pretend
the route has disappeared into certainty. This ending keeps the route visible,
and because the route is visible, the conclusion can be trusted where it has
earned trust.

The book ends by trusting routes, not thunder.

That is the realistic ending: enough record to infer, enough humility to keep
the record open.

The book closes there because that is where science lives.

It lives between refusal and overreach, between the courage to conclude and
the discipline to revise. That middle place is not dramatic, but it is
strong. It is where finite public knowledge becomes real enough to use.

The procession has been a training in that middle place. Marks without
custody were not enough. Residues without discipline were not enough. Staging
without overhead honesty was not enough. Transport without decision was not
enough. Decision gathers them and says: now infer, but remember how you got
here.

That memory is the safeguard. Forget it, and inference becomes slogan. Keep
it, and even a finite conclusion can be worthy of trust, action, and later
correction.

This is the final realism the book can offer. The universe is not handed to us
as a finished object, and public science does not need to pretend that it
is. The universe becomes answerable where records have made it answerable,
and decisions become responsible where those records have earned their
weight. Nothing more is required for seriousness. Nothing less deserves the
name.

So the ending is spare on purpose: not triumph over finitude, but fidelity
inside it.

The book has no better promise to make, and no stronger one it can honestly
keep.

That is why \lean{INFERRED} is an ending, not an apology.

It is the word that lets the route remain true.
\end{movementclose}
